\documentclass[10pt]{article}
\usepackage[margin=20mm]{geometry}
\usepackage{amsmath,amssymb,mathtools,fontspec,unicode-math,xurl,hyperref}
\IfFontExistsTF{DejaVu Serif}{\setmainfont{DejaVu Serif}}{\setmainfont{Latin Modern Roman}}
\setmathfont{latinmodern-math.otf}
\setlength{\parindent}{0pt}\setlength{\parskip}{6pt}\sloppy
\title{Two quantitative corrections to the stated sheeting estimates}
\author{AI-assisted technical review}
\date{}
\begin{document}
\maketitle
Abstract: This scoped review gives explicit stable minimal graphs contradicting two quantitative clauses in the publisher version of Bellettini’s paper. The curvature bounds in Theorems 6 and 7 require an inverse-radius factor. Independently, the normal-deficit square printed in Theorem 7 is fourth order in the slope, so the asserted curvature bound linear in its square-root excess fails even at unit radius. Every geometric hypothesis is verified for the counterexamples. The six-dimensional curvature-iteration endpoint and the tilt Caccioppoli algebra pass the independent checks described here. The findings concern the printed quantitative guarantees; they do not refute the graphical decomposition or the Bernstein conclusions. Scientific adjudication remains pending.\par
\section*{Source identity, conventions, and scope}
The reviewed work is Costante Bellettini, Extensions of Schoen–Simon–Yau and Schoen–Simon theorems via iteration à la De Giorgi, Inventiones mathematicae 240 (2025), 1–34, DOI https://doi.org/10.1007/s00222-025-01317-0 [B]. Only the supplied 34-page publisher PDF is assessed. PDF and printed page numbers coincide. No assertion is made about a later correction or a different version.\par
The complete extracted text was read by the native scientific helper. Independent proof checks concentrate on the scale and amplitude dependence of Theorems 6–7, the curvature iteration at dimension six, the tilt truncation algebra, and the recurrence in Lemma B.1. Eighteen original pages were visually inspected by the helper; the primary executor rechecked the decisive statement and proof images. Full compactness, dimension-reduction, capacity, and cited removable-singularity arguments were not independently reconstructed. Exact expressions, execution receipts, and inspection identities accompany this report.\par
Equation reference: eq:E4. Write \[C_R=B_R^n(0)\times(-R,R),\qquad
 g=\sqrt{1-(\nu\cdot e_{n+1})^2}.\] Theorem 4 (p. 4) uses the scale-invariant quantity \[
 E_4(R)^2=R^{-n}\int_{M\cap C_R}g^2\,d\mathcal H^n.\] It assumes \(n\ge2\) and a properly immersed, two-sided, stable minimal hypersurface in \(C_{2R}\setminus\Sigma\), where \(\Sigma\) is closed and has vanishing \(2\)-capacity. Its smallness assumption is \(E_4(R)^2\le k(n)\) for a positive dimensional constant.\par
Equation reference: eq:E7, eq:printed. Theorem 6 (pp. 6–7) adds \(\Sigma=\varnothing\) and \(\sup_{M\cap C_{2R}}|x_{n+1}|<R/2\). Theorem 7 (p. 7) treats proper embeddings and a permitted singular set, but prints a different excess: \[
 E_7(R)^2=R^{-n}\int_{M\cap C_R}(1-\nu\cdot e_{n+1})^2\,d\mathcal H^n .\] The outer square in [eq:E7] (\#eq:E7) was confirmed in the original page image. Both theorems conclude, with their respective excess, that \[
 \sup_{M\cap C_{R/2}}|A_M|\le c(n)E_M(R).\] The displayed \(|A_M|\) is the ordinary second fundamental form; page 3 explicitly includes \(R^{-1}\) in the analogous Theorem 3 bound.\par
\section*{An explicit admissible family}
Equation reference: eq:graph. For any \(n\ge2\) and \(0<a<1/8\), define on \(B_2^n\) \[
 u_a(x)=\frac{\log\cos(ax_2)-\log\cos(ax_1)}{a},
 \qquad M_a=\{(x,u_a(x)):x\in B_2^n\}.\] For \(R>0\), set \(u_{a,R}(x)=Ru_a(x/R)\) on \(B_{2R}^n\) and \(M_{a,R}=RM_a\). Use the upward normal throughout.\par
\subsection*{Minimality and geometric hypotheses}
The only nonzero first and second derivatives are \[u_1=\tan(ax_1),\quad u_2=-\tan(ax_2),\quad
 u_{11}=a\sec^2(ax_1),\quad u_{22}=-a\sec^2(ax_2).\] The mixed derivative is zero. Consequently the minimal-graph numerator is exactly \[(1+u_2^2)u_{11}+(1+u_1^2)u_{22}
 =a\sec^2(ax_2)\sec^2(ax_1)
  -a\sec^2(ax_1)\sec^2(ax_2)=0.\] Extra coordinates contribute zero. Thus these are minimal hypersurfaces in every indicated dimension, and dilation preserves minimality.\par
The cosine factors are positive on \(B_2^n\), so the graphs are smooth, embedded, and two-sided. They have no boundary inside the open cylinder. They are proper there: a compact subset of \(C_2\) projects to a compact set strictly inside \(B_2^n\), and the graph preimage of its intersection is closed in that compact set. Take \(\Sigma=\varnothing\), which meets every singular-set requirement.\par
For completeness, stability follows from an actual variational argument. The upward normal, extended vertically, has unit length and zero divergence by the minimal-graph equation. Its flux is invariant under compactly supported oriented deformations; the graph attains this flux as its area, while each competitor has area at least its flux. Equivalently, the positive vertical Jacobi field \(w=\nu\cdot e_{n+1}\) gives \[\int_M\bigl(|\nabla_M\phi|^2-|A|^2\phi^2\bigr)
 =\int_M w^2\left|\nabla_M(\phi/w)\right|^2\ge0\] for compactly supported smooth \(\phi\). Small slope is not being substituted for stability.\par
Equation reference: eq:curv. Since \(|\tan t|\le2|t|\) and \(|\log\cos t|\le t^2\) for \(|t|\le1/4\), \[|u_a(x)|\le a(x_1^2+x_2^2)<4a<\tfrac12 .\] Hence \(\sup|u_{a,R}|\le4aR<R/2\), as required in both theorems. The portion \(M_{a,R}\cap C_R\) is exactly the graph over \(B_R^n\). At the origin, \[
 Du_a(0)=0,\qquad D^2u_a(0)=\operatorname{diag}(a,-a,0,\ldots,0),
 \qquad |A_{M_{a,R}}|(0)=\frac{\sqrt2\,a}{R}.\] This point belongs to \(M_{a,R}\cap C_{R/2}\).\par
\subsection*{Exact excess bounds}
Equation reference: eq:bound4, eq:bound7, eq:invariant. Put \(s=|Du_a|^2\), \(W=\sqrt{1+s}\), and let \(\omega_n\) denote the volume of the unit \(n\)-ball. On \(B_1^n\), \[0\le s\le4a^2(x_1^2+x_2^2)\le4a^2|x|^2.\] Since the area element is \(W\,dx\) and \(\nu\cdot e_{n+1}=W^{-1}\), \[\begin{aligned}
 E_4(1)^2&=\int_{B_1^n}\frac{s}{W}\,dx
 \le C_n a^2,& C_n&=\frac{4n\omega_n}{n+2},\\
 E_7(1)^2&=\int_{B_1^n}\frac{s^2}{W(W+1)^2}\,dx
 \le D_n a^4,& D_n&=\frac{4n\omega_n}{n+4}.\end{aligned}\] For the second bound, \(W(W+1)^2\ge4\) is sufficient. Both excesses are positive for \(a>0\) and can meet any prescribed positive smallness threshold. Dilation multiplies area by \(R^n\) while preserving the normal, so \[
 E_i(M_{a,R},R)=E_i(M_a,1),\qquad i=4,7.\] For example, \(C_3=16\pi/5\) and \(D_3=16\pi/7\). The construction therefore does not depend on the separately presented two-dimensional iteration.\par
\section*{The two supported problems}
Problem 1 (Missing inverse-radius factor).  The arbitrary-radius curvature clause of Theorem 6, and the identical clause of Theorem 7, is false as printed in [eq:printed] (\#eq:printed).\par
Fix \(n\ge2\) and choose \(a>0\) so small that the relevant smallness assumption holds. All hypotheses established in Section 2 hold for every \(R>0\). Equations [eq:curv] (\#eq:curv) and [eq:invariant] (\#eq:invariant) would force \[\frac{\sqrt2\,a}{R}\le c(n)E_i(1)\qquad\text{for all }R>0.\] For any proposed finite \(c(n)>0\), choosing \(0<R<\sqrt2a/[c(n)E_i(1)]\) contradicts this inequality. This is a family of admissible counterexamples, not dimensional analysis alone.\par
A scale-compatible formulation is \[R\sup_{M\cap C_{R/2}}|A_M|\le C(n)E_4(R).\] The statement does not restrict \(R\) to one. Allowing \(c\) to depend on \(R\), or silently renormalizing \(|A|\), changes the claim. Theorem 3 on page 3 and the ordinary interior PDE estimates invoked on page 29 support restoring the length scale. The unscaled \(C^{1,\alpha}\) norm of the gradient also needs a convention, but the pointwise curvature clause already establishes the failure without any ambiguity about that norm.\par
Problem 2 (The printed excess has incompatible small-amplitude order).  At \(R=1\), the final curvature bound of Theorem 7 is false when its excess is defined by the printed formula [eq:E7] (\#eq:E7). Restoring \(R^{-1}\) alone does not repair it.\par
For the upward-oriented graph [eq:graph] (\#eq:graph), [eq:curv] (\#eq:curv) and [eq:bound7] (\#eq:bound7) give \[\frac{|A_{M_a}|(0)}{E_7(1)}
 \ge\frac{\sqrt2}{\sqrt{D_n}\,a}\longrightarrow\infty
 \quad(a\downarrow0).\] For any positive constants \(k(n)\) and \(c(n)\), choose \[0<a<\min\left\{\frac18,\,
 \left(\frac{k(n)}{D_n}\right)^{1/4},\,
 \frac{\sqrt2}{c(n)\sqrt{D_n}}\right\}.\] Then \(E_7(1)^2<k(n)\) and all geometric and height assumptions hold, but \(|A|(0)>c(n)E_7(1)\). The normal is fixed upward and tends to \(e_{n+1}\); reversal of orientation plays no role.\par
The quantity used in Theorem 4 and in the tilt iteration is \(1-(\nu\cdot e_{n+1})^2\), whereas the printed Theorem 7 quantity is \((1-\nu\cdot e_{n+1})^2\). The former is quadratic and the latter quartic in a small slope. Replacing the latter by the former, together with restoring the inverse-radius factor, is consistent with the stated proof. This review verifies that these changes remove the two exhibited contradictions. It does not independently prove the entire corrected sheeting theorem. Retaining the printed excess would instead require at least a square-root dependence on \(E_7\) in the small-amplitude regime, with an appropriately proved theorem.\par
\subsection*{Numerical illustration and counterevidence}
The analytic proofs above are decisive. An independently authored script also checked the minimal-graph residual and coefficient identities exactly with SymPy. Actual mpmath quadrature in dimension two gives: \[\begin{array}{c|ccc}
 a&E_4(1)&E_7(1)&|A|(0)/E_7(1)\\ \hline
 0.04&0.0501325691&0.0008186615&69.0988\\
 0.02&0.0250662829&0.0002046653&138.1977\\
 0.01&0.0125331414&0.0000511663&276.3953
\end{array}\] These are numerical illustrations, not interval-certified enclosures or substitutes for the proof.\par
The defects may have short typographical repairs, but each printed theorem clause has a complete counterexample satisfying its hypotheses. Under a standard that rates a demonstrated failure of an asserted quantitative guarantee as Major, the supported assessment is Major, confined to these quantitative clauses. There is no counterexample here to the graphical decomposition, the correctly normalized unit-scale estimate using \(E_4\), Theorems 1–3, or the dimension-six Bernstein conclusion. A whole-paper rejection does not follow from this report.\par
\section*{Independent checks that pass}
\subsection*{The curvature iteration at dimension six}
Pages 11–14 use \(p=2n/(n-2)\), increasing truncation levels \(k_\ell\), and the integrals \(S_\ell=\int_{M\cap B_{R_\ell}}(|A|-k_\ell)_+^p\). The cubic Hölder step has exponents \[\frac{3(n-2)}{2n}\quad\text{on }S_\ell,\qquad
 \frac{6-n}{2n}\quad\text{on the superlevel measure}.\] At \(n=6\) these are \(1\) and \(0\). A positive measure exponent is unnecessary at that endpoint. The other terms and superlevel estimates lead, after choosing the final level \(d=1/R\), to a two-step recurrence with nonlinear power \(n/(n-2)=3/2>1\): \[S_{\ell+1}\le C B^\ell
 R^{-2n(n-4)/(n-2)^2}S_{\ell-1}^{\,n/(n-2)}.\] The normalization \(Y_\ell=R^{p-n}S_\ell\) cancels the radius exponent exactly. Passing to alternate indices yields a standard superlinear recurrence. Thus this calculation supports, rather than invalidates, the six-dimensional endpoint. For \(n=7\) the cubic step would need a negative measure exponent; this argument does not claim that extension.\par
For the recurrence in Lemma B.1 (pp. 32–33), \(x_{\ell+1}\le AB^\ell x_\ell^{1+\alpha}\) with \(B>1\), the threshold \(x_1\le(AB^{1+1/\alpha})^{-1/\alpha}\) is consistent with the decay factor \(B^{-1/\alpha}\). These checks do not certify every explicit numerical overestimate in the paper’s epsilon constants.\par
\subsection*{Tilt coercivity}
For \(g=\sqrt{1-(\nu\cdot e)^2}\), the algebra on pages 24–25 gives \[(g-k)^2-\left(1-\frac1n\right)\left(1-\frac{k}{g}\right)
 +\left(1-\frac{k}{g}\right)(1-g^2)
 =\left(1-\frac{k}{g}\right)\left(\frac1n-kg\right).\] On the active set \(g>k\), with \(0\le g\le1\) and \(0\le k\le1/(2n)\), the final factor is at least \(1/(2n)\). The truncation range therefore supplies the claimed coercivity. The improved factor \(1-1/n\) follows from the trace-free shape-operator inequality \(|Av|^2\le(n-1)|A|^2|v|^2/n\).\par
The tilt recurrence has nonlinear exponent \(1+2/n\), and its radius factor cancels after normalizing \(R^{-n}G_\ell\). Its starting quantity is the \(L^2\) norm of \(g\), namely \(E_4\). This provides an additional proof-level check of the distinction in Problem Problem 2 (\#prob:amplitude).\par
\section*{Disposition and remaining scope}
Item | Recorded conclusion\par
Problem Problem 1 (\#prob:scale) | Counterexample to arbitrary-radius curvature bound; restore length scale.\par
Problem Problem 2 (\#prob:amplitude) | Counterexample at unit radius; printed excess has fourth-order slope dependence.\par
Unaffected conclusions | No refutation of sheeting structure or the Bernstein statements is established.\par
Repair status | Necessary corrections identified; full repaired general theorem not independently certified.\par
Scientific decision | Pending; no human reviews supplied and no independent acceptance completed.\par
The capacity and area-growth extensions on pages 22–25 and the full compactness and dimension-reduction chain remain substantive dependencies outside the independent proof scope. Missing reproofs are not counted as mathematical defects. No author code was executed, and no lack of code is treated as a defect. No claim is made to first discovery or to the absence of a later erratum.\par
\section*{Supporting checks}
The native helper `/root/check\_0057` performed bounded source and proof checks and the independent symbolic and numerical computations. These roles are recorded separately from the controller’s UI model observation and from the runtime self-description; backend identity was not independently authenticated.\par
Rendering alone is never recorded as page inspection. This is an initial scoped review, not independent acceptance.\par
1 C. Bellettini, Extensions of Schoen–Simon–Yau and Schoen–Simon theorems via iteration à la De Giorgi, Invent. Math. 240 (2025), 1–34. Supplied publisher version. https://doi.org/10.1007/s00222-025-01317-0.\par
\end{document}
